\chapter{Resultados Parciais e Discussão}\label{cap:resultados}

    Este capítulo apresenta os resultados parciais obtidos ao longo da primeira fase do trabalho, dedicada à construção e validação dos Hamiltonianos propostos para a representação do TSP no formalismo de computação quântica de variáveis contínuas. Os experimentos foram organizados em torno de duas formulações distintas: a primeira, baseada no mapeamento direto de variáveis binárias para operadores de posição com comportamento binário imposto energeticamente; e a segunda, fundamentada em uma codificação por níveis de Fock, na qual cada qumode representa uma posição da rota e seu nível de ocupação identifica a cidade visitada. Para cada formulação, descrevem-se os algoritmos implementados, os experimentos conduzidos e os resultados obtidos, sempre contrastados com soluções de referência produzidas por busca exaustiva clássica.

    Todas as simulações descritas foram executadas em um servidor institucional da Universidade Federal do ABC(UFABC), denominado Beluga, equipado com processador AMD EPYC 7453 de 56 núcleos e 112 threads, 1 TB de memória RAM e duas GPUs NVIDIA RTX A5500 de 24 GB cada, sob o sistema operacional Ubuntu 24.04.3 LTS. A implementação dos Hamiltonianos e a simulação dos sistemas quânticos foram realizadas com o auxílio da biblioteca QuTiP (\textit{Quantum Toolbox in Python}) \cite{qutip5}, que fornece estruturas para a representação e manipulação de operadores e estados quânticos em espaços de Hilbert de dimensão finita. 

    Os algoritmos utilizados nessa primeira modelagem são descritos a seguir. O \autoref{alg:hamiltoniano_qumode_binario} apresenta a construção do Hamiltoniano qumode-binário para o TSP, definindo os operadores de número associados a cada aresta dirigida e os termos de penalização que impõem as restrições combinatórias do problema. O \autoref{alg:indice_aresta} descreve o mapeamento entre arestas dirigidas e índices de qumodes. Após a obtenção de um estado candidato, o \autoref{alg:indices_para_matriz} converte as ocupações dos modos em uma matriz binária de arestas, enquanto o \autoref{alg:matriz_para_caminho} verifica se essa matriz representa um ciclo Hamiltoniano válido.

    \begin{algorithm}[htbp]
    \caption{Construção do Hamiltoniano qumode-binário para o TSP}
    \label{alg:hamiltoniano_qumode_binario}
    \begin{algorithmic}[1]
    \Require Matriz de distâncias $d_{ij}$, penalidades $\alpha,\beta,\gamma$, dimensão local $\mathrm{dim}$, parâmetro $\mathrm{incluir\_p}$
    \Ensure Hamiltoniano final $H_{\mathrm{final}}$

    \State $N \gets$ número de cidades
    \State $M \gets N(N-1)$ \Comment{um qumode para cada aresta dirigida $i \to j$, com $i \neq j$}

    \For{$k \gets 0$ até $M-1$}
        \State construir o operador de destruição $a_k$
        \State $n_k \gets a_k^\dagger a_k$
    \EndFor

    \State $C \gets 0$
    \For{$i \gets 0$ até $N-1$}
        \For{$j \gets 0$ até $N-1$}
            \If{$i \neq j$}
                \State $k \gets \Call{IndiceAresta}{i,j,N}$
                \State $C \gets C + d_{ij}n_k$
            \EndIf
        \EndFor
    \EndFor

    \State $H_1 \gets 0$
    \For{$k \gets 0$ até $M-1$}
        \State $H_1 \gets H_1 + \alpha\left[n_k(n_k-1)\right]^2$
    \EndFor

    \State $H_2 \gets 0$
    \For{$i \gets 0$ até $N-1$}
        \State $S_{\mathrm{saida}} \gets \sum_{j \neq i} n_{\mathrm{idx}(i,j)}$
        \State $H_2 \gets H_2 + \beta(S_{\mathrm{saida}}-1)^2$
    \EndFor

    \State $H_3 \gets 0$
    \For{$j \gets 0$ até $N-1$}
        \State $S_{\mathrm{entrada}} \gets \sum_{i \neq j} n_{\mathrm{idx}(i,j)}$
        \State $H_3 \gets H_3 + \gamma(S_{\mathrm{entrada}}-1)^2$
    \EndFor

    \State $H_{\mathrm{final}} \gets C + H_1 + H_2 + H_3$

    \If{$\mathrm{incluir\_p}$}
        \State construir o termo cinético $P$
        \State $H_{\mathrm{final}} \gets H_{\mathrm{final}} + P$
    \EndIf

    \State \Return $H_{\mathrm{final}}$
    \end{algorithmic}
    \end{algorithm}

    \begin{algorithm}[htbp]
    \caption{Mapeamento de uma aresta dirigida para o índice do qumode}
    \label{alg:indice_aresta}
    \begin{algorithmic}[1]
    \Require Cidades $i,j$ e número de cidades $N$
    \Ensure Índice $k$ associado à aresta $i \to j$

    \If{$i=j$}
        \State lançar erro
    \EndIf

    \State $k \gets 0$
    \For{$u \gets 0$ até $N-1$}
        \For{$v \gets 0$ até $N-1$}
            \If{$u \neq v$}
                \If{$u=i$ e $v=j$}
                    \State \Return $k$
                \EndIf
                \State $k \gets k+1$
            \EndIf
        \EndFor
    \EndFor
    \end{algorithmic}
    \end{algorithm}


    \begin{algorithm}[htbp]
        \caption{Conversão dos índices de ocupação em matriz binária de arestas}
        \label{alg:indices_para_matriz}
        \begin{algorithmic}[1]
        \Require Tupla de ocupações $\mathrm{indices}$ e número de cidades $N$
        \Ensure Matriz binária $X$ ou $\emptyset$

        \State criar matriz $X \in \{0,1\}^{N \times N}$ inicialmente nula
        \State $k \gets 0$

        \For{$i \gets 0$ até $N-1$}
            \For{$j \gets 0$ até $N-1$}
                \If{$i \neq j$}
                    \State $\mathrm{occ} \gets \mathrm{indices}[k]$

                        \If{$\mathrm{occ} \notin \{0,1\}$}
                            \State \Return $\emptyset$
                        \EndIf

                        \State $X_{ij} \gets \mathrm{occ}$
                        \State $k \gets k+1$
                    \EndIf
                \EndFor
            \EndFor

        \State \Return $X$
        \end{algorithmic}
    \end{algorithm}


    \begin{algorithm}[htbp]
    \caption{Validação da matriz binária como ciclo Hamiltoniano}
    \label{alg:matriz_para_caminho}
    \begin{algorithmic}[1]
    \Require Matriz binária de arestas $X$
    \Ensure Caminho fechado válido ou $\emptyset$

    \State $N \gets$ dimensão de $X$

    \If{alguma linha de $X$ não soma $1$}
        \State \Return $\emptyset$
    \EndIf

    \If{alguma coluna de $X$ não soma $1$}
        \State \Return $\emptyset$
    \EndIf

    \State $\mathrm{caminho} \gets [0]$
    \State $\mathrm{atual} \gets 0$
    \State $\mathrm{visitados} \gets \{0\}$

    \For{$t \gets 1$ até $N-1$}
        \State encontrar $\mathrm{prox}$ tal que $X_{\mathrm{atual},\mathrm{prox}}=1$

        \If{$\mathrm{prox}$ não existe ou não é único}
            \State \Return $\emptyset$
        \EndIf

        \If{$\mathrm{prox} \in \mathrm{visitados}$}
            \State \Return $\emptyset$
        \EndIf

        \State adicionar $\mathrm{prox}$ ao caminho
        \State adicionar $\mathrm{prox}$ a $\mathrm{visitados}$
        \State $\mathrm{atual} \gets \mathrm{prox}$
    \EndFor

    \State encontrar $\mathrm{prox}$ tal que $X_{\mathrm{atual},\mathrm{prox}}=1$

    \If{$\mathrm{prox} \neq 0$}
        \State \Return $\emptyset$
    \EndIf

    \State adicionar $0$ ao final de $\mathrm{caminho}$

    \If{$|\mathrm{visitados}| \neq N$}
        \State \Return $\emptyset$
    \EndIf

    \State \Return $\mathrm{caminho}$
    \end{algorithmic}
    \end{algorithm}




    O algoritmo clássico de força bruta, descrito no \autoref{alg:brute_force}, foi utilizado como referência para validação em ambas as modelagens. Esse algoritmo percorre exaustivamente todas as permutações possíveis das cidades, calcula o custo total de cada rota e retorna o caminho de menor custo. Embora computacionalmente inviável para instâncias de grande porte, sua exatidão o torna ideal para validação em instâncias pequenas, como as utilizadas nos experimentos desenvolvidos.

    
    \begin{algorithm}
        \caption{Algoritmo de Força Bruta}
        \label{alg:brute_force}
        \begin{algorithmic}[1]
        \Require matriz de distâncias \texttt{dist}
        \Ensure melhor custo e melhor caminho
        \State $n \gets \text{len}(\texttt{dist})$
        \State \texttt{vertices $\gets$ [1, \dots, n-1]}
        \State \texttt{best\_cost $\gets$ inf}
        \State \texttt{best\_path $\gets$ None}
        \ForAll{\texttt{perm} em \texttt{permutations(vertices)}}
            \State \texttt{path $\gets$ (0,) + perm + (0,)}
            \State \texttt{cost $\gets$ 0}
            \For{$i \gets 0$ até $\text{len(path)} - 2$}
                \State \texttt{cost $\gets$ cost + dist[path[i]][path[i+1]]}
            \EndFor
            \If{\texttt{cost < best\_cost}}
                \State \texttt{best\_cost $\gets$ cost}
                \State \texttt{best\_path $\gets$ path}
            \EndIf
        \EndFor
        \State \Return \texttt{best\_cost, best\_path}
        \end{algorithmic}
    \end{algorithm}

    Os testes parciais foram conduzidos para $N \in \{3, 4\}$, onde $N$ denota a quantidade de vértices. Para $N = 3$, o algoritmo de força bruta identificou o circuito hamiltoniano de menor custo como $V_0 \rightarrow V_1 \rightarrow V_2 \rightarrow V_0$, resultado compatível com o obtido pelo algoritmo quântico, que produziu o caminho $V_0 \rightarrow V_2 \rightarrow V_1 \rightarrow V_0$, sendo ambos equivalentes a menos de ordem e com custos mínimos iguais. Analogamente, para $N = 4$, o algoritmo de força bruta retornou o caminho $V_0 \rightarrow V_1 \rightarrow V_2 \rightarrow V_3 \rightarrow V_0$, ao passo que o algoritmo quântico produziu $V_0 \rightarrow V_3 \rightarrow V_2 \rightarrow V_1 \rightarrow V_0$, novamente com custos mínimos idênticos. 

    Cabe ressaltar, contudo, que tais experimentos foram restritos a instâncias de pequena dimensão, em razão do elevado custo computacional associado à simulação de sistemas de variáveis contínuas em hardware clássico. Para um sistema composto por $M$ qubits, a representação explícita do vetor de estados requer memória que cresce como $O(2^M)$; já no caso de $M$ qumodes truncados na base de Fock com corte $d$, esse custo escala como $O(d^M)$. Como a dimensão efetiva do espaço de Hilbert depende diretamente de $d$, a simulação clássica de sistemas contínuos torna-se rapidamente proibitiva à medida que $M$ cresce. Nesse contexto, por instâncias de pequena dimensão entendem-se aquelas em que os valores de $M$ e de $d$ permitem armazenar e manipular numericamente os estados e operadores do sistema de forma computacionalmente viável.

    % Como sequência da proposta, foi desenvolvida uma segunda versão do Hamiltoniano, na qual a representação da rota é realizada por meio de variáveis inteiras $n_k \in \{1,2,\dots,N\}$, em que o índice $k$ identifica o vértice $V_k$ e o valor $n_k$ indica a posição em que esse vértice é visitado no percurso. Dessa forma, a configuração do percurso é descrita pelo vetor $(n_1, n_2, \dots, n_N)$, cuja $k$-ésima entrada corresponde ao vértice $V_k$ e cujo valor indica sua ordem de visita na rota. Para exemplificar, considere um estado da forma $\ket{2413}$: o vértice $V_1$ é visitado na segunda posição, $V_2$ na quarta, $V_3$ na primeira e $V_4$ na terceira, resultando na rota $V_3\rightarrow V_1\rightarrow V_4\rightarrow V_2$.

    Como sequência da proposta, foi desenvolvida uma segunda versão do Hamiltoniano, baseada em uma codificação mais natural para o TSP: em vez de associar a cada vértice uma variável que indica sua posição na rota, associa-se a cada \textit{posição} da rota uma variável que indica qual cidade é visitada nessa posição. Formalmente, define-se $c_k \in \{0, 1, \dots, N-1\}$ como a cidade visitada na $k$-ésima posição da rota, de modo que a configuração completa do percurso é descrita pelo vetor $(c_0, c_1, \dots, c_{N-1})$.

    Para exemplificar, considere $N=4$ cidades e o estado $\ket{2\ 0\ 3\ 1}$: a primeira posição da rota é ocupada pela cidade $V_2$, a segunda por $V_0$, a terceira por $V_3$ e a quarta por $V_1$, resultando na rota $V_2 \rightarrow V_0 \rightarrow V_3 \rightarrow V_1 \rightarrow V_2$. Nessa representação, cada qumode $k$ codifica uma posição da rota, e seu nível de Fock $c_k$ identifica a cidade visitada naquela posição. Uma rota válida corresponde, portanto, a um estado em que os valores $c_0, c_1, \dots, c_{N-1}$ formam uma permutação de $\{0, 1, \dots, N-1\}$, garantindo que cada cidade seja visitada exatamente uma vez.

    Nesse contexto, o Hamiltoniano de custo associado ao percurso é dado por:
    $$
    H = \sum_{k=1}^N d_{n_k, n_{k+1}}.
    $$

    Seja $\hat{n}_k$ o operador de número do modo $k$, definido por:
    $$
    \hat{n}_k\ket{a} = a\ket{a},
    $$

    \noindent pode-se escrever o Hamiltoniano correspondente à função objetivo na forma:
    $$
    H = \sum_{k=1}^N \sum_{a=1}^N \sum_{b=1}^N d_{ab}\, \ket{a}\bra{a}_k \otimes \ket{b}\bra{b}_{k+1}.
    $$

    Para impedir a ocorrência de cidades repetidas na rota, introduz-se um termo de penalização definido por:
    $$
    H_{\text{penalty}} = \sum_{k<l} \sum_{a=1}^N \ket{a}\bra{a}_k \otimes \ket{a}\bra{a}_l.
    $$

    Esse termo atribui energia positiva a configurações nas quais uma mesma cidade ocupa múltiplas posições na rota, enquanto estados que correspondem a permutações válidas não são penalizados. Alternativamente, essa penalização pode ser reescrita por meio do operador:
    $$
    \hat{N}_a = \sum_{k=1}^N \ket{a}\bra{a}_k,
    $$

    \noindent resultando na forma equivalente:
    $$
    H_{\text{penalty}} = \sum_{a=1}^N \left(\hat{N}_a - 1 \right)^2.
    $$

    O Hamiltoniano completo é então definido como:
    $$
    H = H_{\text{dist}} + \lambda H_{\text{penalty}},
    $$

    \noindent onde $\lambda > 0$ é um parâmetro de penalização que controla o peso relativo das restrições combinatórias em relação ao custo de deslocamento. Para garantir que todo estado cujos valores $c_0, c_1, \dots, c_{N-1}$ não formem uma permutação válida de $\{0, 1, \dots, N-1\}$ tenha energia estritamente maior do que a melhor rota válida, é suficiente escolher $\lambda$ de modo que a penalização mínima aplicada a qualquer configuração inválida supere o custo máximo de qualquer rota válida. Uma condição conservadora e de fácil verificação é:
    $$
    \lambda > N \cdot d_{\max},
    $$

    \noindent onde $d_{\max} = \max_{i,j} d_{ij}$ é a maior distância entre quaisquer duas cidades. Com essa escolha, qualquer estado que viole a restrição de permutação receberá uma penalização estritamente maior do que o custo total de qualquer rota válida, assegurando que o estado de menor energia corresponda sempre a uma solução viável.

    Os resultados preliminares obtidos com essa modelagem estão sintetizados na \autoref{tab:resultados_segunda_formulacao}, que apresenta, para cada instância testada, o custo mínimo identificado pelo algoritmo de força bruta clássico, a energia mínima identificada pela simulação clássica do Hamiltoniano, o caminho correspondente e o tempo de execução da simulação.

    \begin{table}[!htb]
    \centering
    \caption{Comparação entre os resultados de custo clássicos e energias quânticas para cada quantidade $N$ de cidades e tempo de execução, referente à segunda modelagem do hamiltoniano.}
    \label{tab:resultados_segunda_formulacao}
    \begin{tabular}{|c|c|c|r|}
    \hline
    \textbf{N} & \textbf{Custo} & \textbf{Energia} & \multicolumn{1}{c|}{\textbf{Tempo}} \\ \hline
    3 & 19 & 19,0 & 6 ms \\ \hline
    4 & 23 & 23,0 & 16 ms \\ \hline
    5 & 12 & 12,0 & 86 ms \\ \hline
    6 & 19 & 19,0 & 1 s 229 ms \\ \hline
    7 & 41 & 41,0 & 1 min 7 s \\ \hline
    8 & 26 & 26,0 & 2 h 53 min \\ \hline
    \end{tabular}
    \fonte{Autoria própria (2026)}
    \end{table}

    Em todas as instâncias avaliadas, o caminho de menor custo identificado pelo algoritmo clássico de força bruta coincidiu com o caminho de menor energia obtido pelo algoritmo quântico proposto, confirmando a consistência da formulação. Em cada simulação, a dimensão dos qumodes foi definida como igual ao número de vértices do grafo, adotando corte $d = N$, escolha que assegura expressividade suficiente para representar todas as permutações válidas da rota, embora valores maiores de $d$ possibilitem representação mais fiel ao custo de maior demanda computacional.

    % {\red Senão finalizar até sexta}

    A formulação apresenta uma vantagem estrutural relevante em relação às abordagens baseadas em qubits: um grafo de $N$ vértices pode ser representado com apenas $N$ qumodes com corte $d = N$, enquanto formulações baseadas em qubits, como QUBO ou Ising com codificação \textit{one-hot}, requerem $N^2$ variáveis binárias. Essa diferença se reflete diretamente na dimensão do espaço de estados, conforme ilustrado na \autoref{tab:comparacao}.

    \begin{table}[!htbp]
        \centering
        \caption{Comparação estrutural entre a formulação QUBO \textit{one-hot} 
        e a codificação em qumodes proposta.}
        \label{tab:comparacao}
        \begin{tabular}{lcc}
        \hline
        \textbf{Característica} & \textbf{QUBO \textit{one-hot}} & 
        \textbf{Qumodes (proposta)} \\
        \hline
        Variáveis & $N^2$ binárias & $N$ qumodes \\
        Corte local & --- & $d = N$ \\
        Dimensão do espaço & $2^{N^2}$ & $N^N$ \\
        Exemplo $N=6$ & $2^{36} \approx 6{,}9 \times 10^{10}$ & 
        $6^6 = 46.656$ \\
        Exemplo $N=8$ & $2^{64} \approx 1{,}8 \times 10^{19}$ & 
        $8^8 = 16.777.216$ \\
        Restrições & Termos de penalização & Termos de penalização \\
        Hardware natural & Qubits & Fotônico/CV \\
        \hline
        \end{tabular}
        \fonte{Autoria Própria (2026)}
    \end{table}

    Cabe ressaltar que essa redução no número de variáveis não implica necessariamente vantagem computacional: $N$ qumodes com dimensão local $N$ correspondem a um espaço de dimensão $N^N$, cujo custo de simulação clássica ainda cresce rapidamente com $N$, conforme evidenciado pelos tempos da \autoref{tab:resultados_segunda_formulacao}. A comparação estrutural é, portanto, uma análise de codificação, não de desempenho.

    Os algoritmos utilizados nessa segunda modelagem são descritos a seguir. O \autoref{alg:projector_k_a} define os projetores locais $P_{k,a}$, que identificam a ocupação do nível $a$ no registrador associado à posição $k$ do caminho. Esses projetores são empregados no \autoref{alg:hamiltoniano_tsp_qudit} para construir o Hamiltoniano do problema, composto por um termo de distância $H_{dist}$, responsável por codificar o custo do percurso, e por um termo de penalização $H_{pen}$, que evita a repetição de cidades. Em seguida, o \autoref{alg:validacao_caminho} verifica se uma sequência de índices constitui uma permutação válida das cidades e a converte em um ciclo fechado. Finalmente, o \autoref{alg:analisar_estados_base} percorre os estados da base computacional, calcula suas energias esperadas e identifica a configuração válida de menor energia.

        
\begin{algorithm}[!htbp]
\caption{Construção do projetor local}
\label{alg:projector_k_a}
\begin{algorithmic}[1]
\Require posição $k$, nível $a$, número de registradores $N$, dimensão local $\mathrm{dim}$
\Ensure operador projetor $P_{k,a}$

\State $\mathrm{ops} \gets [\ ]$
\For{$m \gets 0$ até $N-1$}
    \If{$m = k$}
        \State adicionar $|a\rangle\langle a|$ em $\mathrm{ops}$
    \Else
        \State adicionar $\mathbb{I}_{\mathrm{dim}}$ em $\mathrm{ops}$
    \EndIf
\EndFor
\State \Return produto tensorial dos elementos de $\mathrm{ops}$
\end{algorithmic}
\end{algorithm}

\begin{algorithm}[!htbp]
\caption{Construção do Hamiltoniano do TSP em qumodes}
\label{alg:hamiltoniano_tsp_qudit}
\begin{algorithmic}[1]
\Require matriz de distâncias $d_{ij}$, penalização $\lambda$, dimensão local $\mathrm{dim}$
\Ensure Hamiltoniano $H$

\State $N \gets$ número de linhas de $d_{ij}$

\If{$\mathrm{dim}$ não informado}
    \State $\mathrm{dim} \gets N$
\EndIf

\If{$\mathrm{dim} < N$}
    \State lançar erro
\EndIf

\For{$k \gets 0$ até $N-1$}
    \For{$a \gets 0$ até $N-1$}
        \State $P_{k,a} \gets \Call{ProjetorLocal}{k,a,N,\mathrm{dim}}$
    \EndFor
\EndFor

\State $H_{\mathrm{dist}} \gets 0$
\For{$k \gets 0$ até $N-1$}
    \State $k' \gets (k+1)\bmod N$
    \For{$a \gets 0$ até $N-1$}
        \For{$b \gets 0$ até $N-1$}
            \State $H_{\mathrm{dist}} \gets H_{\mathrm{dist}} + d_{ij}[a,b]\,P_{k,a}P_{k',b}$
        \EndFor
    \EndFor
\EndFor

\State $H_{\mathrm{pen}} \gets 0$
\For{$k \gets 0$ até $N-1$}
    \For{$l \gets k+1$ até $N-1$}
        \For{$a \gets 0$ até $N-1$}
            \State $H_{\mathrm{pen}} \gets H_{\mathrm{pen}} + P_{k,a}P_{l,a}$
        \EndFor
    \EndFor
\EndFor

\State \Return $H_{\mathrm{dist}} + \lambda H_{\mathrm{pen}}$
\end{algorithmic}
\end{algorithm}

\begin{algorithm}[!htbp]
\caption{Validação e conversão de índices em caminho}
\label{alg:validacao_caminho}
\begin{algorithmic}[1]
\Require sequência $\mathrm{indices}$, número de cidades $N$
\Ensure caminho válido ou valor nulo

\If{$|\mathrm{indices}| \neq N$}
    \State \Return falso
\EndIf

\If{$\mathrm{sorted}(\mathrm{indices}) \neq [0,1,\dots,N-1]$}
    \State \Return falso
\EndIf

\State \Return tupla $(\mathrm{indices}[0], \ldots, \mathrm{indices}[N-1], \mathrm{indices}[0])$
\end{algorithmic}
\end{algorithm}

\begin{algorithm}[!htbp]
\caption{Análise dos estados da base computacional}
\label{alg:analisar_estados_base}
\begin{algorithmic}[1]
\Require Hamiltoniano $H$, matriz de distâncias $\mathrm{dist}$, número de cidades $N$, dimensão local $\mathrm{dim}$
\Ensure melhor energia, melhor caminho e lista de resultados válidos

\If{$\mathrm{dim}$ não informado}
    \State $\mathrm{dim} \gets N$
\EndIf

\State $\mathrm{melhor\_energia} \gets +\infty$
\State $\mathrm{melhor\_indices} \gets \emptyset$
\State $\mathrm{melhor\_path} \gets \emptyset$
\State $\mathrm{resultados\_validos} \gets [\ ]$

\ForAll{$\mathrm{indices} \in \{0,\dots,\mathrm{dim}-1\}^N$}
    \If{$\mathrm{indices}$ não formam uma permutação válida de $\{0,\dots,N-1\}$}
        \State \textbf{continue}
    \EndIf

    \State construir o estado tensorial correspondente
    \State calcular a energia esperada de $H$
    \State converter $\mathrm{indices}$ em caminho fechado
    \State calcular o custo do caminho
    \State armazenar resultado em $\mathrm{resultados\_validos}$

    \If{$\mathrm{energia} < \mathrm{melhor\_energia}$}
        \State atualizar melhor energia
        \State atualizar melhores índices
        \State atualizar melhor caminho
    \EndIf
\EndFor

\State \Return melhores resultados
\end{algorithmic}
\end{algorithm}

Os resultados apresentados neste capítulo estabelecem a consistência das duas formulações Hamiltonianas propostas: em todas as instâncias avaliadas, os estados de menor energia identificados pelo formalismo de variáveis contínuas coincidiram com os caminhos de menor custo obtidos por busca exaustiva clássica, validando a correção das codificações desenvolvidas. A segunda formulação, em particular, demonstrou uma vantagem estrutural relevante em termos de compacidade da representação, conforme discutido na \autoref{tab:comparacao}, embora o custo de simulação clássica ainda cresça rapidamente com $N$, limitando os experimentos a instâncias de pequena dimensão. Com base nesses resultados, a etapa seguinte do trabalho concentra-se na implementação de algoritmos de otimização quântica variacional, em particular, adaptações de VQE e QAOA ao formalismo de variáveis contínuas, sobre os Hamiltonianos aqui validados, conforme delineado no \autoref{cap:conclusao}.